\begin{corollary}[Finite existence over an arbitrary Weil coefficient field]
\label{c:finite-weil-criterion}
\status{Proved under $C^+_H(X)$; effectiveness requires the stated coefficient data}
Let $X/k$ be smooth projective and connected over an arbitrary
algebraically closed field and of positive dimension, let $H$ be a
Weil cohomology with coefficient field $K$ of characteristic zero,
and let $o$ be a
rational zero-cycle of degree one.  Suppose $C^+_H(X)$, meaning
that a rational algebraic correspondence realizes parity on
$H^*(X)$.  The finite criterion (i)--(iii) of
Theorem~\ref{c:finite-super-criterion} still exists, with a finite
list of \emph{rational} marked diagrams, and is equivalent to
injectivity of
\[
 T_k^{\CH}(X,o)\longrightarrow H^*(X^k)
 \quad\hbox{for every }k.
\]
Here injectivity refers to rational linear relations, as required
for cycles in $\CH_\Q$.  No rational form of the whole cohomology
algebra is assumed.  The finite list is effectively constructible
if a finite presentation of the $\Q$-algebra generated by its
chosen structure constants is supplied.  For arbitrary elements
of an unspecified coefficient field, only finite existence is
asserted.  In dimension zero, $X=\operatorname{Spec}(k)$ and
all the requested relations hold immediately.
\end{corollary}

\begin{proof}
Choose a homogeneous $K$-basis of $H^*(X)$ and let
$A_\Q=\Q^{a|b}$ be a formal supervector space of the same
dimensions.  It is not asserted to be a rational form of the cup
algebra.  Build $G_\Q$, $W_\Q$ and $R_\Q=\mathcal O(W_\Q)$
from this formal space exactly as above.  The actual structure
tensors give a $K$-point $t$ of $W_\Q$.  Replace the orbit ideal by
\[
 I_\Q=\ker\left(R_\Q\longrightarrow
        \mathcal O(G_\Q)\otimes_\Q K\right),
 \qquad f\longmapsto\bigl(g\longmapsto f(g\cdot t)\bigr).
\]
The map is a $G_\Q$-comodule map, with the regular action on the
first target factor and trivial action on the coefficient field.
Its kernel is consequently a rational $G_\Q$-stable ideal.
Equivalently, the kernel assertion follows by tensoring its
underlying exact sequence with $\mathcal O(G_\Q)$ and using the
coaction identity.  A rational covariant $F$ satisfies $F(t)=0$
if and only if its coordinates lie in $I_\Q$: in one direction
equivariance gives $F(g\cdot t)=gF(t)$ on all superalgebra-valued
points, and in the other we evaluate at the identity.

Noetherianity of $R_\Q$ gives finite ideal generators; closing
their span in a bounded polynomial degree gives a finite rational
stable $E\subset I_\Q$.  Lemma~\ref{c:super-projective-tensors}
and every subsequent projective lift take place over $\Q$.
Evaluation now takes place in $K$-cohomology.  The trace cancellation
still gives the positive integer $\dim_\Q P_0$, since
$\operatorname{End}(\mathbf1)=\Q$ embeds in $K$.
The same necessity argument uses rational permutation graphs and
the Frobenius normalization in Chow.  Thus the finite iff survives
unchanged.

For the effectiveness assertion, let $c_1,\ldots,c_m\in K$ be
all the coordinates of $t$, including the inverse pairing, and
write $B=\Q[c_1,\ldots,c_m]$.  If a finite presentation
$B=\Q[z_1,\ldots,z_m]/J$ is given, the orbit map lands already
in $\mathcal O(G_\Q)\otimes B$, which injects into
$\mathcal O(G_\Q)\otimes K$.  Orbit-graph elimination including
the equations of $J$ therefore computes $I_\Q$.  Without a
presentation of $B$, no such computational claim follows merely
from the abstract knowledge of the elements $c_i$.

If instead one asks for every $K$-linear relation to vanish in
$\CH_\Q\otimes_\Q K$, the same proof is performed entirely over
$K$, using the full $K$-orbit ideal.  That is a separate
coefficient convention; it is not deduced by tensoring an
injective rational cycle map with $K$ and identifying its target
with the original cohomology space.
\end{proof}

\begin{corollary}[The sign hypothesis over an algebraic closure of a finite field]
\label{c:finite-field-sign}
\status{Proved for the stated cohomology theories}
For $X/\overline{\F}_p$, the sign hypothesis in
Corollary~\ref{c:finite-weil-criterion} holds for $\ell$-adic
cohomology, $\ell\ne p$, and crystalline cohomology with
characteristic-zero coefficients.  More generally it holds for
every Weil theory satisfying the Poincar\'e duality, weak
Lefschetz and zeta-function hypotheses in section~II of
\citet{c_katzmessing1974}.  Thus in this setting the finite
criterion has no additional sign-conjecture hypothesis.
\end{corollary}

\begin{proof}
Descend $X$ to a smooth projective geometrically connected variety
over a finite field $\F_q$.  The finitely many points and rational
coefficients representing $o$ also descend after enlarging
$\F_q$; the parity construction itself does not depend on this
choice of $o$.  Put
$Q_i(T)=\det(T-F\mid H^i(X))$.
\citet[Theorem~1, pp.~74--75; Theorem~2(1), pp.~76--77]{c_katzmessing1974}
proves that $Q_i\in\Q[T]$ and that all its roots have absolute
value $q^{i/2}$ in every complex embedding, with the same
polynomials for the cohomology theories just specified.
Consequently the nonconstant $Q_i$ are pairwise coprime.
By the Chinese remainder theorem there are polynomials
$r_i\in\Q[T]$ with $r_i\equiv1\pmod{Q_i}$ and
$r_i\equiv0\pmod{Q_j}$ for $j\ne i$.
Cayley--Hamilton shows that $r_i(F)$ acts as the projector onto
$H^i(X)$.  Each is a rational combination of graphs of iterates
of the Frobenius morphism, and hence an algebraic correspondence.
Their sum over even $i$ gives the required sign projector after
base extension to $\overline{\F}_p$.  This is the explicit
polynomial argument in Theorem~2(1), so it does not require
Frobenius semisimplicity.
\end{proof}
